\scanpage{14}
\[
(18)\qquad\widehat f:\overset{\deleted{\unsure{\wedge}}}{X}\longrightarrow Y
\]
se factorise en
\diagram{\begin{tikzcd}[ampersand replacement=\&]
(19)\quad X\arrow[r,"\widehat f_0","\sim"']\&Y'\arrow[r,"u"]\&Y
\end{tikzcd}}
où $u$ est un $C_0$-morphisme, et $\widehat f_0$ un $\widehat C_0$-\uline{isomorphisme}
(\uline{NB} $Y'$ est un objet $C_0$-isom. à $X$), et où $u$ est
un \unsure{morphisme|épimorphisme} de $C_0$. Dans le cas où on
suppose que pour toute flèche $u:Y'\mathbin{\unsure{-|\longrightarrow}}Y$ dans $C$,
l'hom. composé
\diagram{\begin{tikzcd}[ampersand replacement=\&]
\operatorname{Aut}_Y(Y')\arrow[r]\&\operatorname{Aut}(Y')\arrow[r,no head]\&\operatorname{Aut}(Y')^{\wedge}
\end{tikzcd}}
\uline{injectif}, on voit que $Y'$ est déterminé
en termes de $\widehat f$ à $Y$-isom. unique près. \deleted{\unsure{De}} \unsure{La|Le}
\inserted{recette} \deleted{\unsure{plus, pour}} pour composer deux morphismes définis
ainsi, $\widehat f=u\widehat f_0$ et $\widehat g=v\widehat g_0$
\margin{\uline{NB} Le foncteur \unsure{can.|canonique} $C_0\longrightarrow\widehat C_0$ est conservatif (i.e. $\unsure{\widehat f_0|f_0}$ est $C_0$-iso $\Longleftrightarrow$ c'est un $\widehat C_0$-iso) \deleted{\unsure{De}} plus, Tout $\widehat C_0$-isom. est limite de $C_0$-isom.\ldots\ de façon précise,
$\operatorname{Isom}_{\widehat C_0}(X,Y)\simeq\operatorname{Isom}_{C_0}(X,Y)\mathbin{\wedge}_{\operatorname{Aut}_{C_0}(X)}\widehat{\operatorname{Aut}}_{C_0}(X)$
et $\operatorname{Aut}_{\widehat C_0}(X)\simeq\widehat{\operatorname{Aut}}_{C_0}(X)$.}
\diagram{\begin{tikzcd}[ampersand replacement=\&]
(20)\quad Y\arrow[r,no head,"\deleted{\unsure{g}}\widehat g_0","\sim"']\&Z'\arrow[r,no head,"v"]\&Z,
\end{tikzcd}}
et les mettre sous la $=$ forme standard,
est le suivant : on approche $\widehat g_0:Y\overset{\sim}{\longrightarrow}Z'$
(qui est un $\widehat C_0$-iso) par
un \deleted{\unsure{vrai}} $C_0$-iso \inserted{morphisme} \deleted{$\unsure{u}$} $w$
de telle façon que l'on ait
\[
\widehat g_0=w\widehat\varphi\qquad\text{avec}\quad\widehat\varphi\in\widehat H_u\subset\widehat{\operatorname{Aut}}_{C_0}(Y)
\]
\diagram{\begin{tikzcd}[ampersand replacement=\&,column sep=small]
\&\&Y'\arrow[r,"u"]\&Y\arrow[dr,"w"]\&\&\\
(21)\quad X\arrow[r,"\widehat f_0","\sim"']\&Y'\arrow[r,"u"]\arrow[ur,"\widehat\psi?" ]\&Y\arrow[rr,"\widehat g_0","\sim"']\arrow[ur,"\widehat\varphi?"]\&\&Z'\arrow[r,"v"]\&Z
\end{tikzcd}}
\margin{\uline{NB} Il faut savoir interpréter en termes de $C_0$ \unsure{déjà} la relation $\widehat\varphi u=u\widehat\psi$ \unsure{comme|comme un} \unsure{signifiant|signe signifiant} $(\widehat\psi,\widehat\varphi)\in\widehat G_u$, où $G_u\subset\operatorname{Aut}(X)\times\operatorname{Aut}(Y)$ est le stabilisateur de $u$.}
(ce qui est possible, puisque $\widehat H_u$ est ouvert dans $\widehat{\operatorname{Aut}}_{C_0}(Y)$),
et on choisit $\widehat\psi\in\widehat{\operatorname{Aut}}(Y')$ tel qu'on ait $\widehat\varphi u=u\widehat\psi$
donc on aura
\[
(22)\qquad\widehat g\widehat f=(v\widehat g_0)(u\widehat f_0)
=(v\deleted{\unsure{g}}w\widehat\varphi)(u\widehat f_0)
=\underbrace{(vwu)}_{\in\mathrm{Fl}\,C_0}\underbrace{(\widehat\psi\widehat f_0)}_{\in\mathrm{Flisom.}(\widehat C_0)}
\]
\margin{13}
